\documentclass[a4paper]{article}
\usepackage{amsfonts}
\usepackage{amsmath}
\usepackage{amssymb}
\usepackage{graphicx}
\usepackage{extarrows}

\newcommand{\integral}{\displaystyle\int}

\begin{document}

\noindent \large Auteur: Seda Jusjajeva\\
\noindent \large Studierichting: Informatica\\
\noindent \large Oefeningenreeks 4A nr. 58.8\\

\medskip

\normalsize

$\integral \dfrac{1}{4\cos x + 3\sin x}dx$ \\

t-formules: $t = \tan \dfrac{x}{2} \Rightarrow dt = \dfrac{1}{2}\left(1 + \tan^2 \dfrac{x}{2}\right)dx \Leftrightarrow dx = \dfrac{2}{t^2 + 1}dt$ \\

$= \integral \dfrac{1}{4\dfrac{1 - t^2}{1 + t^2} + 3\dfrac{2t}{1 + t^2}}\cdot\dfrac{2}{t^2 + 1}dt$ \\

$= \integral \dfrac{1 + t^2}{4 - 4t^2 + 6t}\cdot\dfrac{2}{t^2 + 1}dt$ \\

$= 2\integral \dfrac{1}{4 - 4t^2 + 6t}dt$ \\

$= \integral \dfrac{1}{-2t^2 + 3t + 2}dt$ \\

$-2t^2 + 3t + 2 = 0 \Leftrightarrow (t - 2)(-2t - 1) = 0$ \\

$\Rightarrow\dfrac{1}{-2t^2 + 3t + 2} = \dfrac{A}{t - 2} + \dfrac{B}{-2t - 1}$ \\

$A = \left.\left(\dfrac{1}{-2t - 1}\right)\right|_{x = 2} = -\dfrac{1}{5} $ \\

$B = \left.\left(\dfrac{1}{t - 2}\right)\right|_{x = -\tfrac{1}{2}} = -\dfrac{2}{5} $ \\

$\Rightarrow \integral\dfrac{1}{-2t^2 + 3t + 2}dt = -\dfrac{1}{5}\integral\dfrac{1}{t - 2}dt - \dfrac{2}{5}\integral\dfrac{1}{-2t - 1}dt$ \\

$= -\dfrac{1}{5}\ln |t - 2| + \dfrac{2}{10} \ln |-2t - 1| + c$ \\

$= -\dfrac{1}{5}\ln \left|\tan \dfrac{x}{2} - 2\right| + \dfrac{1}{5} \ln \left|-2\tan \dfrac{x}{2} - 1\right| + c$ \\



\begin{thebibliography}{999}
%\bibitem{a} Referentie
\end{thebibliography}
\end{document}